\paragraph{Normalized Anchoring Index (\nai).}
For each item $i$ with low-anchor samples $\{y^L_j\}$ and
high-anchor samples $\{y^H_j\}$, we define
\begin{equation}
  \text{NAI}_i = \frac{\bar{y}^H_i - \bar{y}^L_i}{60},
  \label{eq:nai}
\end{equation}
where $\bar{y}$ denotes the sample mean and 60 is the anchor gap.
$\text{NAI} = 1$ indicates complete susceptibility;
$\text{NAI} = 0$ indicates no anchoring.

\paragraph{Distributional metrics.}
We complement \nai with per-item Total Variation Distance (\tvd)
and Wasserstein-1 distance (\wone) between the low-anchor and
high-anchor empirical distributions.

\paragraph{Relative Reduction (\rr).}
For a mitigation $m$ applied to model $k$, we compute
\begin{equation}
  \text{RR}_{k,m,i} = 1 - \frac{|\text{NAI}^m_i|}{|\text{NAI}^{B0}_i|},
  \label{eq:rr}
\end{equation}
and report the mean \rr with bootstrap 95\% CIs and
Benjamini--Hochberg adjusted $p$-values.

\paragraph{Controlled-history slope ($\beta_{\text{norm}}$).}
For the conversation-history suite, we regress the model's
stage-2 estimate on the controlled anchor
$a \in \{10, 30, 50, 70, 90\}$ and normalize the OLS slope:
$\beta_{\text{norm}} = \beta \cdot 60 / 80$.
